\documentclass[12pt]{article}
\usepackage{amsmath, amsthm, amssymb}
\usepackage[utf8]{inputenc}
\usepackage[margin=1in]{geometry}

\newtheorem{theorem}{Theorem}

\begin{document}
\section{Intervals in Dense Orders}
\begin{theorem}
In an o-minimal structure $\mathcal{M} = (M, <, \dots)$, for any $a, b \in M$ with $a < b$, the open interval $(a, b)$ is infinite.
\end{theorem}
\begin{proof}
By the definition of an o-minimal structure, the underlying order $<$ is dense. This means that for any $x, y \in M$ with $x < y$, there exists $z \in M$ such that $x < z < y$.

Let $x_0 = b$. Since $a < x_0$, the density property guarantees there exists $x_1$ such that $a < x_1 < x_0$. 
Assume we have recursively chosen $x_1, \dots, x_k$ such that:
\[ a < x_k < \dots < x_1 < b \]

Because $a < x_k$, density guarantees an $x_{k+1}$ such that $a < x_{k+1} < x_k$. This constructs an infinite, strictly decreasing sequence within $(a, b)$. Since the sequence is strictly decreasing, all elements are distinct. Thus, the interval is infinite.
\end{proof}



\end{document}